\section{A fixed-window test and the localization question}
\label{sec:fixed-window}

The first realization test can be made before the first arithmetic
delay. For $0<L<\log2$, no nontrivial integer translation acts on
$I_L$. The local term and both rational pole terms remain. On this
window we can state the complete target, calculate a normalized
Wilson readout, and exclude that readout within a precise regularity
class. We can also formulate a localization test for the required
cancellations. The detailed calculations and proofs are in
\cite[Sections~S10--S12]{Supplement}.

\subsection{The local and pole terms constrain the response}

Let $A_L=-G_{0,L}$ be the required first shift derivative. Away
from the diagonal its causal kernel is
\begin{equation}\label{eq:main-required-response}
 b(u)=2\frac{e^{-u/2}}{1-e^{-2u}}-4\cosh(u/2)
     =2\frac{e^{-5u/2}}{1-e^{-2u}}-2e^{u/2},
 \qquad 0<u<L.
\end{equation}
The first pole cancels the lowest gamma mode. The other leaves
the signed growing exponential in the second expression.
The kernel behaves as $1/u$ at zero; it must be interpreted
together with a contact subtraction. In the hard-delay-cutoff
convention, on $C_c^\infty(I_L)$,
\begin{equation}\label{eq:main-contact}
 (A_Lf)(x)=\lim_{\epsilon\downarrow0}
 \left[c_\epsilon f(x)+
 \int_\epsilon^{x+L/2}b(u)f(x-u)\dd u\right],
 \qquad
 c_\epsilon=\log\epsilon+\gamma+\log(2\pi)+O(\epsilon).
\end{equation}
An integral whose upper limit is at most $\epsilon$ means zero.
This finite part is fixed by $w_0$ in the stated cutoff scheme,
not a new normalization choice. Its derivation, including the
zero-extension tail, is in \cite[Section~S10.1]{Supplement}.

At positive shift, absorb the first rational factor into the
gamma ratio. The factor whose compression is the full arithmetic
transfer on this window becomes
\begin{equation}\label{eq:main-small-factor}
 K_\omega^{<}(p)=
 \pi^\omega
 \frac{\Gamma((p+5/2-\omega)/2)}
      {\Gamma((p+5/2+\omega)/2)}
 \frac{p-1/2-\omega}{p-1/2+\omega}.
\end{equation}
This is an arithmetic identity, not yet a field-theory
construction. Its locally integrable positive-shift kernel has
a singularly concentrating identity limit; there is no separate
delta channel at positive shift
\cite[Section~S10.2]{Supplement}.

\subsection{A normalized causal Wilson readout that fails the test}

Use the ordinary reflected endpoint pairing, with the physical
polarizations $q_2,\bar q_2$, a common reference color fiber and
the adjoint scalar map $N=-MR$. For the smooth polynomial bulges
of the companion, deform a positive height by
$\eta(\omega)=\eta_0+\omega$, $\eta_0>0$, keeping the endpoints
fixed. Let $\mathsf K_\omega(x,y)$ be the radially weighted
reflected correlator. A definite linear response is
\begin{equation}\label{eq:main-readout}
 (\mathcal V_{\omega,L}f)(x)=f(x)+
 \int_{-L/2}^{x}
 [\mathsf K_\omega(x,y)-\mathsf K_0(x,y)]f(y)\dd y.
\end{equation}
It has exact identity initial data and causal support. The
retarded integration in logarithmic radius is part of the
specified readout, not a consequence of Euclidean reflection.
It avoids inverting the compact free Gram operator.

In the independent abelian Gaussian probe model defined in
\cite[Sections~S10.3--S10.4]{Supplement}, the free endpoint term
has zero variation, while the Wilson dressing gives a nonzero
variation at positive bulge height. The unequal-radius
calculation and a uniform bound show that
$\partial_\omega\mathcal V_{0,L}$ is nevertheless a bounded
integral operator. It therefore cannot equal $-G_{0,L}$ on the
smooth core: for $f_N(x)=e^{\ii Nx}\chi(x)$ with nonzero
$\chi\in C_c^\infty(I_L)$,
\[
 \frac{Q_{0,L}[f_N]}{\norm\chi^2}
 =\log N-\log(2\pi)+o(1).
\]
No bounded derivative, even with a bounded local multiplier,
has this quadratic-form growth. At positive shift there is a
second obstruction: \eqref{eq:main-readout} is identity plus a
compact operator, whereas $V_{\omega,L}$ is compact.
The proofs and regulator scope are in
\cite[Section~S10.5]{Supplement}.

These exclusions do not cover all Wilson constructions or a
controlled singular limit of regulated responses. They show
what must change in the tested observable-to-operator map.
Localizing the same regular readout would not remove the
obstruction.

\subsection{Localization as a test of the surviving modes}

Localization is a candidate mechanism for deriving the
cancellations in \eqref{eq:main-required-response}. Pestun's
fluctuation calculation organizes cancellations by an
equivariant complex~\cite{Pestun}. Baker's open-line discussion
already points toward a lower-dimensional defect theory
\cite[Section~3.3]{Baker}. Wang's defect-localization framework
contains endpoint bilocals joined by gauge transport
\cite[Section~5.1, equation~(5.5)]{WangDefects}; its unevaluated
normal-fluctuation determinant remains an explicit assumption
\cite[Section~4.4]{WangDefects}.

An auxiliary calculation makes the proposed test concrete. Put
$a_n=2n+1/2$ and retain one complex bosonic Gaussian for each
$n=1,\ldots,N$, together with one complex Grassmann pair:
\begin{equation}\label{eq:main-bf-product}
 Z_N(p)=(p-1/2)\prod_{n=1}^N(p+a_n)^{-1}.
\end{equation}
Then exact finite-product algebra gives
\begin{equation}\label{eq:main-bf-ratio}
 \pi^\omega\lim_{N\to\infty}N^{-\omega}
 \frac{Z_N(p-\omega)}{Z_N(p+\omega)}
 =K_\omega^{<}(p).
\end{equation}
Removing the lowest boson accounts for the first pole
cancellation, and the Grassmann factor supplies the remaining
rational ratio. These modes and the normalization have been
chosen from the arithmetic target. Their presence is not a
localization theorem or even a specified supersymmetric action.

The physical test is whether a defect-compatible supercharge,
boundary conditions and observable produce these surviving
contributions and the required finite normalization.
An especially sharp comparison, independent of an additive
$p$-independent generator constant, is
\begin{equation}\label{eq:main-mode-test}
 \partial_p a_0^{<}(p)
 =2\sum_{n=1}^{\infty}\frac1{(p+2n+1/2)^2}
  -\frac2{(p-1/2)^2}.
\end{equation}
The construction and its missing physical dictionary are
spelled out in \cite[Section~S11]{Supplement}. A localization
deformation coefficient multiplying a $Q$-exact term cannot
simply be identified with the physical shift. Likewise a
partition ratio does not yet act on arbitrary boundary inputs.
The original conformal semicircle remains a relevant control,
outside the restricted constant-Poincare-charge endpoint
obstruction.

\subsection{A predictive boundary calculation and its spatial limitation}
\label{subsec:main-localized-control}

A specified free three-dimensional $\mathcal N=4$ hypermultiplet
on a hemisphere provides a concrete localization control
\cite{DPYHiggs,DedushenkoGluing}. With its Higgs boundary
polarization, put $\epsilon=(4\pi r)^{-1}$ and
$x=|Y|^2/\epsilon$. The vacuum is proportional to $e^{-x}$,
whose boundary Fourier--Mellin amplitude is
$\Gamma(s)/\sqrt{2\pi}$ at $s=1/2-\ii\sigma$.
Separated scalar endpoints with identity transport give
$\epsilon x e^{-x}$ and hence
$\epsilon\Gamma(s+1)/\sqrt{2\pi}$. The gamma factor and this
insertion follow from the selected physical problem; they do not
require the arithmetic tower to be prescribed. Their derivation,
ordering and ordinary-adjoint checks are in
\cite[Sections~S12.1--S12.2]{Supplement}.

The transform acts on a boundary field's magnitude. It does not
yet transform spatial logarithmic radius. In an adapted flat frame,
a protected scalar is
$O(t)=q_1(t)+tq_2(t)/(2r)$, with
$\mathcal Qq_2(t)=\chi(t)$ a nonzero free fermion.
Ordinary spatial dilation gives
\begin{equation}\label{eq:main-localized-nonclosure}
 \mathcal Q[D,O(t)]=-\frac{t}{2r}\chi(t).
\end{equation}
Thus the fixed protected representation does not receive the
ordinary radial evolution on arbitrary closed representatives.
The compensating R rotation restores closure but gives the exact
twisted dilation, trivial on the tested local classes. The centered
hemisphere geometry and the scope of this obstruction are in
\cite[Section~S12.3]{Supplement}.

There is a distinct canonical-Hodge alternative. Selecting the
kernel of $D-R_H$ first permits a restricted radial semigroup.
For an elementary scalar, however, a spatial harmonic has
$D=\ell+1/2$ and $R_H=1/2$, so only $\ell=0$ survives.
The projected endpoint response is $e^{-u/2}$.
The omitted even-angular channel is
\[
 \frac{e^{-5u/2}}{1-e^{-2u}},
\]
precisely the positive tower appearing in
\eqref{eq:main-required-response}. Composite protected states can
have a matching counting character after additional selections,
but an elementary endpoint does not excite that character.
A trace, a boundary amplitude and an ordinary-adjoint response
are different tests \cite[Sections~S12.4--S12.5]{Supplement}.

These calculations provide a physical precedent for a gamma
amplitude and a specific obstruction to its spatial interpretation
in this free model. They do not rule out an enlarged or interacting
sector. Any continuation must specify how it retains the missing
modes and their endpoint matrix elements, before fitting a
spectral-parameter map or finite normalization.
A polynomial-degree comparison using an unproved affine map is
conditional; changing the offset can change that degree
\cite[Section~S12.6]{Supplement}.

Finally, signed mode cancellation does not establish an
ordinary positive norm. The independent cumulative-storage
identity of Section~\ref{sec:storage} and the spatial gluing
estimates of Section~\ref{sec:path} remain separate requirements.
Neither this test nor the auxiliary product supplies an
arithmetic Loewner driver.
